\section{Empirical Evaluation}
\label{sec:evaluation}

This section quantifies the overhead $\Mtype\langle T \rangle$ adds over
$\Utype\langle T \rangle$ at the operator and sampling levels, and the cost
of the SPRT-gated lifting operator across the three regimes of the
hypothesis test. All benchmarks are shipped with the crate and can be
reproduced; the figures reported here are taken from a representative run
on the development host.

\subsection{Methodology}
\label{sec:eval:methodology}

The crate ships two Criterion~\cite{criterion} benchmark suites:
\texttt{uncertain\_benchmarks} for the existing $\Utype\langle T \rangle$
type and \texttt{maybe\_uncertain\_benchmarks} for
$\Mtype\langle T \rangle$. The two suites share their structure so that
benchmark functions are pairwise comparable on the per-operation level.
Each benchmark function reports the mean wall-clock time per iteration
together with a $95\%$ confidence interval estimated by Criterion's
internal bootstrap.

Benchmarks were collected on macOS (Darwin 25.5) with the crate compiled
in release mode using the default Cargo bench profile. Each Criterion run
draws $100$ measurements per benchmark and reports the median and the $95\%$
confidence band. The hardware configuration used to produce the numbers
reported in this paper is recorded with the benchmark artifact and should
be consulted for absolute timing reference. We refrain from cross-machine
comparisons; the numbers here are relative cost ratios, which are stable
across hardware in our spot checks.

\subsection{Sampling overhead}
\label{sec:eval:sampling}

\Cref{tab:sampling} compares per-sample cost between $\Utype\langle f64
\rangle$ and $\Mtype\langle f64 \rangle$ across six paired benchmarks. The
left column gives the baseline cost of sampling a $\Utype$ value; the
middle column gives the corresponding cost for a $\Mtype$ value built from
the same underlying distribution. The right column is the ratio.

\begin{table}[h]
\centering
\small
\begin{tabular}{@{}lrrr@{}}
\toprule
Operation & $\Utype\langle f64 \rangle$ & $\Mtype\langle f64 \rangle$ & Ratio \\
\midrule
Point sample                       & $237$~ns & $510$~ns & $2.15\times$ \\
Normal sample / \textsc{fromUncertain} & $286$~ns & $496$~ns & $1.73\times$ \\
Bernoulli-gated sample             & $246$~ns & $400$~ns & $1.63\times$ \\
Arithmetic chain $(a+b)\cdot c$    & $384$~ns & $747$~ns & $1.95\times$ \\
Add and sample                     & $477$~ns & $762$~ns & $1.60\times$ \\
Complex chain and sample           & $593$~ns & $718$~ns & $1.21\times$ \\
\bottomrule
\end{tabular}
\caption{Per-sample latency for paired benchmarks on $\Utype\langle f64
\rangle$ and $\Mtype\langle f64 \rangle$. The presence axis adds a constant
factor between $1.2\times$ and $2.2\times$ to per-sample cost. The
``Bernoulli-gated sample'' row pairs \texttt{sampling\_bernoulli} on the
underlying $\Utype\langle \Bool \rangle$ with
\texttt{sampling\_maybe\_bernoulli} on a $\Mtype\langle f64 \rangle$ value
constructed via \textsc{fromBernoulli}.}
\label{tab:sampling}
\end{table}

The pattern is consistent. The presence axis costs one extra sample plus
one extra cache lookup; the structural overhead does not grow with the
complexity of the value-side graph. The ``complex chain and sample''
benchmark, which combines three Bernoulli-gated readings with arithmetic
and a division, shows the smallest relative penalty ($1.21\times$) because
the value-side work dominates the constant cost of the presence draw.

A second observation matters for embedded callers. The
$\textsc{alwaysNone}$ constructor is the cheapest case: sampling an
always-absent $\Mtype\langle f64 \rangle$ takes $266$~ns, which sits
slightly above the cost of sampling a single Bernoulli point distribution
($246$~ns). The short-circuit in the sampler is empirically observable;
the value branch is not evaluated.

\subsection{Graph-construction overhead}
\label{sec:eval:construction}

\Cref{tab:construction} reports the cost of building a single operator
node, with no sampling. Construction times measure the work done by the
operator overload to extend the computation graph and produce a new owning
value.

\begin{table}[h]
\centering
\small
\begin{tabular}{@{}lrr@{}}
\toprule
Operator & $\Utype\langle f64 \rangle$ & $\Mtype\langle f64 \rangle$ \\
\midrule
Binary \texttt{+} (graph only)    & $14.7$~ns & $42.9$~ns \\
Binary \texttt{*} (graph only)    & ---       & $44.4$~ns \\
Unary \texttt{-} (graph only)     & ---       & $14.7$~ns \\
Map (graph only)                  & $26.5$~ns & ---       \\
\bottomrule
\end{tabular}
\caption{Graph-construction latency for selected operators. The
$\Mtype\langle f64 \rangle$ binary operators cost roughly $2.9\times$ a
$\Utype\langle f64 \rangle$ binary operator because each
$\Mtype$ binary op extends two computation graphs. Unary negation on
$\Mtype$ takes the same time as the value-only operation because the
presence axis is left untouched by the operator algebra
(\Cref{eq:neg}).}
\label{tab:construction}
\end{table}

The $\Mtype$-binary-to-$\Utype$-binary ratio of $2.9\times$ has a direct
explanation: each binary operation on $\Mtype$ constructs both an
\texttt{ArithmeticOp} node in the value graph and a \texttt{LogicalOp::And}
node in the presence graph. The unary-negation row gives an empirical
confirmation of the algebra: unary negation does not touch the presence
axis (\Cref{eq:neg}), and the measured construction cost matches the
$\Utype$ baseline to within measurement noise.

\subsection{Cost of the SPRT-gated lifting operator}
\label{sec:eval:sprt}

The behaviour of \texttt{lift\_to\_uncertain} is dominated by the
underlying SPRT. We benchmark three regimes, each invoking the operator
with threshold $p_\tau = 0.5$, confidence $1 - \alpha = 0.95$, indifference
half-width $\varepsilon = 0.05$, and sample cap $N = 1000$. The presence
probability $p$ of the operand varies across the regimes.

\begin{table}[h]
\centering
\small
\begin{tabular}{@{}lcccc@{}}
\toprule
Regime & Operand $p$ & Threshold $p_\tau$ & Time & Decision \\
\midrule
Confidently present     & $0.95$ & $0.5$ & $448$~ns  & Accept $H_1$ \\
Confidently absent      & $0.10$ & $0.5$ & $470$~ns  & Accept $H_0$ \\
Close to threshold      & $0.55$ & $0.5$ & $5\,803$~ns & Fallback at $N$ \\
\bottomrule
\end{tabular}
\caption{Cost of \texttt{lift\_to\_uncertain} across three SPRT regimes.
The two decisive regimes terminate in a few hundred nanoseconds; the
ambiguous regime exhausts the sample cap and falls back to an empirical
comparison.}
\label{tab:sprt}
\end{table}

The two decisive regimes terminate within the first batch of ten samples
in most invocations, because the log-likelihood ratio after a single
batch already crosses the relevant Wald boundary. Their cost is
proportional to the cost of ten Bernoulli samples on the presence graph
plus a small fixed overhead for boundary computation.

The ambiguous regime exhausts the sample cap. With $p = 0.55$ sitting
precisely on the upper edge of the indifference region $[0.45, 0.55]$, the
SPRT cannot distinguish $H_0$ from $H_1$ within the requested error budget,
and the loop runs all $N = 1000$ samples before falling back to the
empirical comparison $\bar{p} > p_\tau$. The measured ratio of
$\approx 13\times$ between ambiguous and decisive regimes matches this
expectation: roughly $1000/10 = 100\times$ more sampling work modulated
by per-sample cost differences and constant overhead.

\subsection{Discussion}
\label{sec:eval:discussion}

Three takeaways follow from these numbers.

\paragraph{Constant-factor overhead.} The presence axis adds a constant
factor between $1.2\times$ and $2.2\times$ to per-sample cost and a
constant factor of roughly $2.9\times$ to binary graph construction. The
overhead does not grow with the complexity of the underlying computation:
in the ``complex chain'' benchmark the relative penalty is smallest
because the value-side work dominates.

\paragraph{SPRT cost is bimodal.} The two decisive regimes pay only the
cost of a few Bernoulli samples; the ambiguous regime pays a worst case
proportional to $N$. The bimodal behaviour gives callers a clean operating
discipline: choose $\varepsilon$ and $N$ together to bound the worst case,
and expect the typical case to be cheap when the operand's presence
probability lives well away from the threshold. The empirical ratio
quoted above ($\approx 13\times$) is one specific instance; the ratio
scales linearly with the chosen $N$.

\paragraph{When the overhead matters.} For interactive workloads and the
worked examples in~\Cref{sec:examples}, the absolute overhead is a few
hundred nanoseconds per operation. For embedded or hot-loop callers, the
overhead is non-trivial and the choice between $\Mtype$ and $\Utype$
should be deliberate: use $\Mtype$ where sparse data is genuinely a
feature of the input and the cost is justified by the bug class it
prevents; use $\Utype$ where every input is guaranteed to be present and
the presence axis would only add noise.
